% !TeX program = pdflatex
% Biên dịch 2 lần bằng pdfLaTeX, XeLaTeX hoặc LuaLaTeX.
\documentclass[12pt,a4paper,oneside]{report}

\usepackage[a4paper,top=2.5cm,bottom=2.5cm,left=3.2cm,right=2.2cm]{geometry}
\usepackage{iftex}
\ifPDFTeX
  % Nhánh tương thích trình biên dịch mặc định pdfLaTeX trên Overleaf.
  \usepackage[utf8]{inputenc}
  \usepackage[T5]{fontenc}
  \usepackage[vietnamese]{babel}
\else
  % XeLaTeX/LuaLaTeX dùng Unicode và font OpenType trực tiếp.
  \usepackage{fontspec}
  \defaultfontfeatures{Ligatures=TeX}
  \setmainfont{TeX Gyre Termes}
  \setsansfont{TeX Gyre Heros}
  \setmonofont{Latin Modern Mono}[Scale=0.90]
\fi

\usepackage{microtype}
\usepackage{setspace}
\usepackage{graphicx}
\usepackage{booktabs}
\usepackage{longtable}
\usepackage{tabularx}
\usepackage{array}
\usepackage{multirow}
\usepackage{enumitem}
\usepackage{xcolor}
\usepackage{fancyhdr}
\usepackage{titlesec}
\usepackage{tikz}
\usetikzlibrary{arrows.meta,positioning,shapes.geometric,fit}
\usepackage{hyperref}
\usepackage{bookmark}
\usepackage{caption}
\usepackage{float}

% -----------------------------------------------------------------------------
% Thông tin sinh viên -- điền trước khi nộp
% -----------------------------------------------------------------------------
\newcommand{\StudentName}{[HỌ VÀ TÊN SINH VIÊN]}
\newcommand{\StudentID}{[MÃ SỐ SINH VIÊN]}
\newcommand{\ClassName}{[LỚP]}
\newcommand{\AdvisorName}{[GIẢNG VIÊN HƯỚNG DẪN]}
\newcommand{\ReportDate}{Tháng 9 năm 2026}

% -----------------------------------------------------------------------------
% Định dạng tài liệu
% -----------------------------------------------------------------------------
\definecolor{DataOpsBlue}{HTML}{164E63}
\definecolor{DataOpsCyan}{HTML}{0891B2}
\definecolor{DataOpsGreen}{HTML}{15803D}
\definecolor{DataOpsAmber}{HTML}{B45309}
\definecolor{DataOpsRed}{HTML}{B91C1C}
\definecolor{DataOpsGray}{HTML}{475569}
\definecolor{DataOpsLight}{HTML}{F1F5F9}

\hypersetup{
  unicode=true,
  colorlinks=true,
  linkcolor=DataOpsBlue,
  urlcolor=DataOpsCyan,
  citecolor=DataOpsGreen,
  pdftitle={Báo cáo tiến độ tiểu luận chuyên ngành -- DataOps Control Plane},
  pdfauthor={\StudentName}
}

\setstretch{1.3}
\setlength{\parindent}{1.2cm}
\setlength{\parskip}{0.25em}
\setlist[itemize]{leftmargin=1.8em,itemsep=0.25em,topsep=0.4em}
\setlist[enumerate]{leftmargin=2em,itemsep=0.25em,topsep=0.4em}
\setlength{\emergencystretch}{3em}

\newcolumntype{Y}{>{\raggedright\arraybackslash}X}
\newcolumntype{C}[1]{>{\centering\arraybackslash}p{#1}}
\newcolumntype{L}[1]{>{\raggedright\arraybackslash}p{#1}}

\titleformat{\chapter}[display]
  {\normalfont\bfseries\color{DataOpsBlue}}
  {\filleft\Large Chương \thechapter}
  {0.5ex}
  {\titlerule\vspace{1ex}\Huge}
\titleformat{\section}{\Large\bfseries\color{DataOpsBlue}}{\thesection}{0.7em}{}
\titleformat{\subsection}{\large\bfseries\color{DataOpsGray}}{\thesubsection}{0.7em}{}

\pagestyle{fancy}
\fancyhf{}
\fancyhead[L]{\small Báo cáo tiến độ tiểu luận chuyên ngành}
\fancyhead[R]{\small DataOps Control Plane}
\fancyfoot[C]{\thepage}
\renewcommand{\headrulewidth}{0.4pt}
\renewcommand{\footrulewidth}{0pt}

\renewcommand{\contentsname}{Mục lục}
\renewcommand{\listfigurename}{Danh mục hình}
\renewcommand{\listtablename}{Danh mục bảng}
\renewcommand{\figurename}{Hình}
\renewcommand{\tablename}{Bảng}
\renewcommand{\bibname}{Tài liệu và minh chứng tham khảo}

\newcommand{\statusdone}{\textcolor{DataOpsGreen}{\textbf{Đã hoàn thành}}}
\newcommand{\statuspartial}{\textcolor{DataOpsAmber}{\textbf{Một phần}}}
\newcommand{\statusnext}{\textcolor{DataOpsRed}{\textbf{Chưa hoàn thành}}}
\newcommand{\code}[1]{\texttt{#1}}

\begin{document}

% -----------------------------------------------------------------------------
% Trang bìa
% -----------------------------------------------------------------------------
\begin{titlepage}
  \begin{center}
    {\large\bfseries TRƯỜNG ĐẠI HỌC SƯ PHẠM KỸ THUẬT\par}
    {\large\bfseries THÀNH PHỐ HỒ CHÍ MINH\par}
    \vspace{0.35cm}
    {\large\bfseries KHOA CÔNG NGHỆ THÔNG TIN\par}

    \vspace{1.4cm}
    {\color{DataOpsBlue}\rule{0.72\textwidth}{1.2pt}\par}
    \vspace{0.8cm}
    {\LARGE\bfseries BÁO CÁO TIẾN ĐỘ\par}
    \vspace{0.25cm}
    {\Large\bfseries TIỂU LUẬN CHUYÊN NGÀNH\par}
    \vspace{0.8cm}
    {\Large\bfseries\color{DataOpsBlue}
      XÂY DỰNG NỀN TẢNG DATAOPS ĐA NỀN TẢNG\par
      TÍCH HỢP AI AGENT VÀ HYBRID RAG\par
      ĐỂ PHÂN TÍCH, PHỤC HỒI SỰ CỐ PIPELINE DỮ LIỆU\par}
    \vspace{0.8cm}
    {\color{DataOpsBlue}\rule{0.72\textwidth}{1.2pt}\par}

    \vfill
    \begin{tabular}{L{5.0cm} L{8.0cm}}
      \textbf{Sinh viên thực hiện:} & \StudentName \\
      \textbf{Mã số sinh viên:} & \StudentID \\
      \textbf{Lớp:} & \ClassName \\
      \textbf{Giảng viên hướng dẫn:} & \AdvisorName \\
      \textbf{Mã nguồn:} & \url{https://github.com/AndyAnh174/dataops-control-plane} \\
    \end{tabular}
    \vfill
    {\large Thành phố Hồ Chí Minh, \ReportDate\par}
  \end{center}
\end{titlepage}

\pagenumbering{roman}

\chapter*{Thông tin về báo cáo}
\addcontentsline{toc}{chapter}{Thông tin về báo cáo}

Báo cáo này mô tả trạng thái thực tế của dự án tại ngày 20/09/2026. Nội dung được tổng hợp
từ mã nguồn, tài liệu thiết kế, kết quả kiểm thử tự động, bản triển khai demo và hai lần chạy
pipeline end-to-end. Các năng lực đã có mã nguồn nhưng chưa vượt qua kiểm thử thực tế được đánh
dấu là \emph{một phần}, không được xem là hoàn thành ở mức production.

Các thông tin đặt trong dấu ngoặc vuông ở trang bìa cần được sinh viên điền trước khi nộp.
Mọi token, mật khẩu và thông tin bí mật đều bị loại khỏi báo cáo.

\chapter*{Tóm tắt}
\addcontentsline{toc}{chapter}{Tóm tắt}

Đề tài xây dựng một DataOps Control Plane độc lập với nhà cung cấp CI/CD. Hệ thống nhận event,
log và báo cáo chất lượng dữ liệu từ một DataOps Agent chạy trong pipeline; tạo incident khi có
lỗi; thu thập bằng chứng; truy xuất tri thức bằng Hybrid RAG; dùng LLM để đề xuất phân tích nguyên
nhân gốc; sau đó đưa đề xuất qua Policy Engine, cổng phê duyệt và quy trình verification trước khi
đóng incident.

Đến thời điểm báo cáo, các milestone M1--M7 đã có nền tảng mã nguồn: ingestion, incident,
evidence, Data Quality report, Hybrid Retrieval, LangGraph RCA, policy/recovery/verification,
authentication, workspace, project, token và Web UI. Hệ thống đã được đóng gói container, triển
khai public qua HTTPS và chạy thành công một pipeline bình thường cùng một pipeline lỗi có kiểm
soát. Bộ test hiện tại đạt 96 test; bài scan image không còn lỗ hổng mức High/Critical.

Giới hạn quan trọng nhất là LLM RCA chưa ổn định: cả \code{gemma4:e2b} và
\code{qwen3.5:4b} đều trả output không vượt qua schema validation trên incident thực. Do đó,
phần AI hiện đạt mức thử nghiệm, chưa được quyền tự thực hiện recovery. Ngoài ra, Docker Hub,
provider thứ hai, bộ benchmark nhiều fault scenario và UAT đầy đủ theo ba vai trò vẫn là công
việc tiếp theo. Đánh giá tổng thể hiện tại là \textbf{đạt mức demo kỹ thuật, chưa production-ready
cho AI RCA}.

\vspace{0.5em}
\noindent\textbf{Từ khóa:} DataOps, CI/CD, FastAPI, Hybrid RAG, LangGraph, Ollama,
Elasticsearch, Data Quality, RCA, Human-in-the-loop.

\tableofcontents
\listoffigures
\listoftables
\clearpage
\pagenumbering{arabic}

% =============================================================================
\chapter{Tổng quan đề tài}

\section{Bối cảnh và bài toán}

Pipeline dữ liệu thường được ghép từ nhiều công cụ: hệ thống quản lý mã nguồn, CI/CD, engine xử
lý dữ liệu, kho log và hạ tầng triển khai. Khi pipeline thất bại, kỹ sư phải tự tìm run tương ứng,
đọc lượng lớn log, đối chiếu commit, kiểm tra báo cáo chất lượng dữ liệu và tra cứu runbook. Quy
trình này mất thời gian, khó tái lập và phụ thuộc kinh nghiệm cá nhân.

Đề tài giải quyết bài toán trên bằng một control plane đứng ngoài CI/CD. Control plane chuẩn hóa
dữ liệu từ provider, theo dõi vòng đời run, lưu bằng chứng và hỗ trợ xử lý incident. AI chỉ đóng
vai trò phân tích và đề xuất; quyết định thực thi được kiểm soát bởi policy xác định trước.

\section{Mục tiêu}

Các mục tiêu chính của đề tài gồm:

\begin{enumerate}
  \item Chuẩn hóa cách theo dõi pipeline giữa nhiều nền tảng CI/CD.
  \item Thu thập event, log, Data Quality report và metadata theo cùng một định danh run.
  \item Tạo incident và bằng chứng có checksum, nguồn trích dẫn và cơ chế chống trùng.
  \item Kết hợp tìm kiếm từ khóa và vector để truy xuất runbook/incident tương tự.
  \item Tạo RCA có cấu trúc, có citation và có thể kiểm tra bằng máy.
  \item Thực hiện recovery an toàn qua Policy Engine, approval, audit và verification.
  \item Cung cấp Web UI self-hosted để quản lý workspace, project, token, run và incident.
  \item Chứng minh kiến trúc có thể mở rộng sang GitLab CI hoặc Jenkins mà không viết lại core.
\end{enumerate}

\section{Phạm vi hiện tại}

MVP chọn GitHub Actions làm provider đầu tiên và pipeline batch làm loại workload chính. Hệ thống
không thay thế GitHub Actions, không tự merge code, không trực tiếp sửa dữ liệu production và
không gọi LLM ở happy path. Provider thứ hai và đánh giá production-scale nằm ngoài phạm vi bản
demo hiện tại.

Sản phẩm gồm hai phần tách biệt:

\begin{itemize}
  \item \textbf{DataOps Agent:} chạy trong CI/CD, đọc \code{dataops.yaml}, thực thi tuần tự các
  stage, giữ nguyên exit code và gửi telemetry về Control Plane.
  \item \textbf{DataOps Platform/Control Plane:} một deployable FastAPI phục vụ JSON API và Web
  UI; PostgreSQL lưu trạng thái, Elasticsearch lưu log/knowledge, Ollama phục vụ embedding và RCA.
\end{itemize}

% =============================================================================
\chapter{Phân tích yêu cầu và nguyên tắc thiết kế}

\section{Yêu cầu chức năng}

\begin{longtable}{L{1.3cm} L{6.0cm} L{6.5cm}}
\caption{Các nhóm yêu cầu chức năng chính}\label{tab:functional-requirements}\\
\toprule
\textbf{Mã} & \textbf{Yêu cầu} & \textbf{Kết quả mong đợi} \\
\midrule
\endfirsthead
\toprule
\textbf{Mã} & \textbf{Yêu cầu} & \textbf{Kết quả mong đợi} \\
\midrule
\endhead
FR-01 & Đăng ký workspace, project và repository & Mỗi project có cấu hình provider và token tích hợp riêng. \\
FR-02 & Nhận event/log/report & Dữ liệu được validate, correlation theo run và chống ghi trùng. \\
FR-03 & Quản lý incident & Run lỗi tạo tối đa một incident và có state transition rõ ràng. \\
FR-04 & Thu thập evidence & Metadata, log lỗi, commit diff và Data Quality report có citation/checksum. \\
FR-05 & Hybrid Retrieval & BM25 và dense vector được hợp nhất bằng Reciprocal Rank Fusion. \\
FR-06 & Agentic RCA & RCA theo JSON Schema, claim gắn evidence/knowledge ID hợp lệ. \\
FR-07 & Recovery an toàn & Policy quyết định deny/approval; executor idempotent; có verification. \\
FR-08 & Web UI và RBAC & OWNER, OPERATOR, VIEWER có quyền phù hợp; không có public signup. \\
FR-09 & Self-hosted & Có container image, Compose, healthcheck, backup và rollback guide. \\
FR-10 & Đa nền tảng & Core không phụ thuộc provider; thêm adapter mà không đổi domain chính. \\
\bottomrule
\end{longtable}

\section{Yêu cầu phi chức năng}

\begin{itemize}
  \item \textbf{An toàn:} LLM không nhận credential; log được redact; action thay đổi trạng thái
  phải qua policy và approval.
  \item \textbf{Tin cậy:} event, report và recovery có idempotency; incident chỉ đóng khi
  verification pass.
  \item \textbf{Khả năng vận hành:} healthcheck, immutable image tag, backup trước rollout và
  rollback rõ ràng.
  \item \textbf{Khả năng mở rộng:} provider adapter, modular monolith và storage tách theo loại
  dữ liệu.
  \item \textbf{Hiệu năng:} happy path không chờ LLM; chỉ failure path mới chạy evidence/RAG/RCA.
  \item \textbf{Khả năng kiểm toán:} lưu actor, quyết định policy, approval, attempt và kết quả
  verification.
\end{itemize}

\section{Nguyên tắc human-in-the-loop}

Trong bản hiện tại, LLM không có quyền gọi trực tiếp provider hoặc hạ tầng. Model chỉ tạo RCA và
recommended action. Policy Engine là thành phần deterministic kiểm tra môi trường, loại incident,
độ tin cậy và action. Các thao tác \code{RETRY}, \code{QUARANTINE}, \code{ROLLBACK} hoặc
\code{CREATE\_PR} được chuyển qua approval gate trước khi Recovery Executor chạy. Cách thiết kế
này tránh việc một output sai schema hoặc hallucination của model tác động trực tiếp lên production.

% =============================================================================
\chapter{Kiến trúc và luồng hoạt động}

\section{Kiến trúc tổng thể}

\begin{figure}[H]
\centering
\resizebox{\textwidth}{!}{%
\begin{tikzpicture}[
  node distance=12mm and 14mm,
  box/.style={draw=DataOpsBlue, rounded corners=2mm, very thick, align=center,
              minimum width=30mm, minimum height=12mm, fill=DataOpsLight},
  store/.style={draw=DataOpsGray, cylinder, shape border rotate=90, aspect=0.28,
                very thick, align=center, minimum height=13mm, fill=white},
  arrow/.style={-{Latex[length=2.5mm]}, thick, draw=DataOpsGray}
]
  \node[box] (dev) {Developer\\commit/push};
  \node[box, right=of dev] (ci) {GitHub Actions\\Self-hosted runner};
  \node[box, right=of ci] (agent) {DataOps Agent\\\code{dataops.yaml}};
  \node[box, right=of agent, minimum width=38mm] (api) {FastAPI Control Plane\\API + Web UI};
  \node[store, above right=8mm and 16mm of api] (pg) {PostgreSQL};
  \node[store, right=16mm of api] (es) {Elasticsearch};
  \node[box, below right=8mm and 16mm of api] (ollama) {Ollama\\Qwen + BGE-M3};
  \node[box, below=17mm of api] (operator) {Operator\\browser/session};

  \draw[arrow] (dev) -- (ci);
  \draw[arrow] (ci) -- (agent);
  \draw[arrow] (agent) -- node[above, font=\small] {event/log/report} (api);
  \draw[arrow] (api) -- (pg);
  \draw[arrow] (api) -- (es);
  \draw[arrow] (api) -- (ollama);
  \draw[arrow] (operator) -- node[right, font=\small] {HTTPS} (api);
  \draw[arrow, bend right=18] (api.south west) to node[below, font=\small] {recovery dispatch} (ci.south east);
\end{tikzpicture}%
}
\caption{Kiến trúc triển khai hiện tại của DataOps Control Plane}
\label{fig:architecture}
\end{figure}

Core được thiết kế provider-neutral. Mọi thao tác đọc diff, trigger workflow hoặc verification đều
đi qua provider adapter. FastAPI được chọn theo hướng modular monolith: API và HTML/CSS/JavaScript
cùng một deployable, nhưng domain, service và adapter được tách module để có thể mở rộng sau này.

\section{Luồng từ commit đến kết quả}

\begin{enumerate}
  \item Developer commit và push mã nguồn lên GitHub.
  \item GitHub Actions cấp runner, checkout code và gọi \code{dataops-agent@v0}.
  \item Agent đọc \code{dataops.yaml}, thực thi stage, gửi event, log và report bằng project token.
  \item Control Plane cập nhật \code{PipelineRun}; PostgreSQL lưu metadata, Elasticsearch lưu log.
  \item Nếu thành công, run chuyển \code{SUCCESS}; không gọi LLM và không tạo incident.
  \item Nếu thất bại, Control Plane tạo incident, thu thập evidence và truy xuất tri thức.
  \item LangGraph gọi Ollama để tạo RCA có cấu trúc; server kiểm tra schema và citation.
  \item Policy Engine đánh giá action; operator phê duyệt khi policy yêu cầu.
  \item Recovery Executor gọi provider adapter để retry/quarantine/rollback.
  \item Verification callback quyết định incident được \code{RESOLVED} hay cần can thiệp.
\end{enumerate}

\section{Luồng an toàn khi AI sai}

\begin{figure}[H]
\centering
\begin{tikzpicture}[
  node distance=8mm,
  step/.style={draw=DataOpsBlue, rounded corners, thick, align=center,
               minimum width=105mm, minimum height=9mm, fill=DataOpsLight},
  bad/.style={draw=DataOpsRed, rounded corners, thick, align=center,
              minimum width=105mm, minimum height=9mm, fill=red!5},
  arrow/.style={-{Latex[length=2.5mm]}, thick, draw=DataOpsGray}
]
  \node[step] (a) {Evidence Collector + Hybrid Retrieval};
  \node[step, below=of a] (b) {LLM tạo RCA JSON};
  \node[step, below=of b] (c) {Pydantic/schema/citation validation};
  \node[bad, below=of c] (d) {Không hợp lệ: không lưu RCA, không tạo recovery plan tự động};
  \node[step, below=of d] (e) {Chuyển ACTION\_REQUIRED / ANALYSIS\_FAILED và lưu audit};
  \draw[arrow] (a) -- (b);
  \draw[arrow] (b) -- (c);
  \draw[arrow] (c) -- node[right, font=\small, text=DataOpsRed] {fail} (d);
  \draw[arrow] (d) -- (e);
\end{tikzpicture}
\caption{Fail-safe flow đề xuất khi output RCA không hợp lệ}
\label{fig:ai-fail-safe}
\end{figure}

Trạng thái \code{ANALYSIS\_FAILED} trong Hình~\ref{fig:ai-fail-safe} là công việc P0 tiếp theo.
Phiên bản đang chạy hiện giữ incident ở \code{ANALYZING} khi API trả 422; đây là một giới hạn đã
được ghi nhận, không phải trạng thái cuối mong muốn.

% =============================================================================
\chapter{Thiết kế và hiện thực hệ thống}

\section{Công nghệ sử dụng}

\begin{table}[H]
\centering
\caption{Tech stack của phiên bản hiện tại}
\label{tab:tech-stack}
\begin{tabularx}{\textwidth}{L{3.7cm} Y}
\toprule
\textbf{Nhóm} & \textbf{Công nghệ và vai trò} \\
\midrule
Backend/Web & Python 3.11+, FastAPI, Jinja2, HTML, CSS, JavaScript \\
Validation/ORM & Pydantic Settings, SQLModel/SQLAlchemy \\
Metadata & PostgreSQL \\
Log và knowledge & Elasticsearch; Kibana cho quan sát vận hành \\
AI orchestration & LangGraph, Ollama, JSON Schema \\
Embedding & \code{bge-m3:567m}, vector 1024 chiều \\
RCA model thử nghiệm & \code{qwen3.5:4b}; trước đó \code{gemma4:e2b} \\
CI/CD & GitHub Actions, self-hosted runner, \code{dataops-agent@v0} \\
Container/Deploy & Docker, Docker Compose, GHCR, Nginx, Certbot/TLS \\
Quality/Security & pytest, Ruff, Trivy \\
\bottomrule
\end{tabularx}
\end{table}

\section{Các module đã hiện thực}

Mã nguồn tổ chức thành bốn lớp chính:

\begin{itemize}
  \item \code{api/routes}: event, run, log, incident, retrieval, recovery, web authentication,
  member, project, token, onboarding và UI.
  \item \code{services}: pipeline event/log, evidence, GitHub adapter, Elasticsearch, Ollama,
  retrieval, RCA, policy, execution, audit và web identity.
  \item \code{domain}: model quan hệ và state của hệ thống.
  \item \code{web}: template Jinja2 cùng static CSS/JavaScript.
\end{itemize}

Tại thời điểm thống kê, repository có 41 file Python phía ứng dụng với khoảng 5.854 dòng, và
22 file test với khoảng 4.357 dòng. Các con số này chỉ dùng mô tả quy mô hiện tại, không được xem
là chỉ số chất lượng độc lập.

\section{Identity, workspace và token}

Hệ thống không có public signup. Lần khởi động đầu tiên cho phép bootstrap một OWNER. OWNER đăng
nhập bằng server-side session, tạo tài khoản thành viên và gán một trong ba vai trò:
\code{OWNER}, \code{OPERATOR}, \code{VIEWER}. Hệ thống chặn owner tự xóa hoặc tự hạ quyền của
chính mình.

Mỗi project có integration token riêng. Secret chỉ hiển thị một lần, server chỉ lưu hash/prefix;
token có scope và có thể revoke. OWNER có thể xóa project sau khi nhập đúng repository để xác nhận.
Token mất hiệu lực ngay nhưng lịch sử run/incident vẫn được giữ để audit.

\section{Data Quality, evidence và retrieval}

DataOps Agent upload report versioned gồm contract, scenario, kết quả tổng và từng check. Các
scenario hỗ trợ gồm schema drift, missing value, duplicate, range và volume. Khi incident được
tạo, Evidence Collector gom pipeline metadata, failed-stage log, GitHub commit diff và Data
Quality report. Nội dung được giới hạn kích thước, redact secret và chống trùng bằng SHA-256.

Knowledge index tách khỏi raw log. BM25 tìm exact term như tên cột, error code; dense vector tìm
runbook hoặc incident tương tự về ngữ nghĩa. Hai danh sách được hợp nhất bằng Reciprocal Rank
Fusion. Raw log không bị embedding từng dòng nhằm giảm chi phí và tránh nhiễu.

\section{Recovery và verification}

Policy Engine trả một trong ba quyết định: \code{DENIED}, \code{REQUIRE\_APPROVAL} hoặc
\code{AUTO\_APPROVED}. Bản demo cấu hình thận trọng, yêu cầu con người phê duyệt action thay đổi
trạng thái. Recovery attempt có idempotency key, audit trail và số lần retry giới hạn. Việc
dispatch workflow chỉ chứng minh lệnh đã được gửi; incident chỉ đóng sau verification \code{PASSED}.

% =============================================================================
\chapter{Tiến độ thực hiện}

\section{Trạng thái theo milestone}

\begin{longtable}{C{1.2cm} L{4.2cm} L{6.0cm} L{2.4cm}}
\caption{Tiến độ M1--M8 tại ngày 20/09/2026}\label{tab:milestones}\\
\toprule
\textbf{Mốc} & \textbf{Phạm vi} & \textbf{Kết quả hiện tại} & \textbf{Trạng thái} \\
\midrule
\endfirsthead
\toprule
\textbf{Mốc} & \textbf{Phạm vi} & \textbf{Kết quả hiện tại} & \textbf{Trạng thái} \\
\midrule
\endhead
M1 & Ingestion và incident & Event/log API, correlation, idempotency, run state, incident khi failed & \statusdone \\
M2 & Evidence & Metadata, failed log, GitHub diff, checksum, redact, duplicate detection & \statusdone \\
M3 & Data Quality & Report contract, upload trước failed event, các scenario mẫu & \statusdone \\
M4 & Hybrid Retrieval & BM25 + BGE-M3 vector + filter + RRF, knowledge index riêng & \statusdone \\
M5 & Agentic RCA & LangGraph và schema validation đã có; model thực tế chưa sinh output hợp lệ ổn định & \statuspartial \\
M6 & Policy và recovery & Deterministic policy, approval, executor, audit, verification callback & \statusdone \\
M7 & Web Platform & Bootstrap/login, workspace/member/role, project/token, UI, onboarding, guarded delete & \statusdone \\
M8 & Đa nền tảng & Core và adapter contract đã chuẩn bị; chưa tích hợp GitLab/Jenkins thực tế & \statusnext \\
\bottomrule
\end{longtable}

\section{Đối chiếu Definition of Done}

\begin{longtable}{L{7.5cm} L{4.5cm} L{2.0cm}}
\caption{Mức đáp ứng tiêu chí hoàn thành MVP}\label{tab:dod}\\
\toprule
\textbf{Tiêu chí} & \textbf{Minh chứng hiện tại} & \textbf{Mức đạt} \\
\midrule
\endfirsthead
\toprule
\textbf{Tiêu chí} & \textbf{Minh chứng hiện tại} & \textbf{Mức đạt} \\
\midrule
\endhead
Onboard repository GitHub bằng template & Web UI sinh workflow, \code{dataops.yaml}, danh sách Secrets & Đạt \\
Success và failure xuất hiện đúng project & Hai pipeline thực tế đã correlation đúng & Đạt \\
Lưu run, incident, evidence và audit & PostgreSQL/Elasticsearch có dữ liệu; restart không mất volume & Đạt \\
RCA hợp lệ trên incident thực & Cả Gemma và Qwen còn fail schema & Một phần \\
Không có action nguy hiểm vượt policy & AI chỉ đề xuất; approval và executor tách biệt & Đạt \\
Demo retry/quarantine/rollback end-to-end & API/adapter đã có; chưa báo cáo đủ ba kịch bản lặp lại & Một phần \\
Ít nhất năm fault scenario có ground truth & Đã thiết kế F01--F06; mới chạy công khai \code{null\_rate} & Chưa đạt \\
Token không truy cập chéo/revoke tức thời & Scope project, hash token và test auth đã có & Đạt \\
Phân phối self-hosted hoàn chỉnh & GHCR đa kiến trúc có; Docker Hub chưa cấu hình & Một phần \\
Provider thứ hai chứng minh portability & Chưa triển khai GitLab CI/Jenkins & Chưa đạt \\
\bottomrule
\end{longtable}

Với cách chấm theo phạm vi kỹ thuật, dự án hiện hoàn thành khoảng \textbf{80\% khối lượng MVP}.
Đây là ước lượng quản lý tiến độ, không phải số đo production readiness. Phần còn lại có độ rủi ro
cao hơn tỷ lệ khối lượng vì tập trung ở độ tin cậy AI, benchmark, provider thứ hai và phát hành.

% =============================================================================
\chapter{Triển khai, kiểm thử và kết quả demo}

\section{Môi trường triển khai}

Control Plane được triển khai bằng Docker Compose, public qua Nginx và HTTPS tại:

\begin{itemize}
  \item Web Console: \url{https://dataops-console.andyanh.id.vn}
  \item Ứng dụng demo: \url{https://dataops.andyanh.id.vn}
  \item Mã nguồn: \url{https://github.com/AndyAnh174/dataops-control-plane}
\end{itemize}

Release được pin bằng image bất biến:
\code{ghcr.io/andyanh174/dataops-control-plane:sha-5575c6e31bdc62c65bbba3a07bd0bf62f9f24577}.
Trước rollout, PostgreSQL được backup; Compose config được validate và chỉ API container bị
recreate. PostgreSQL, Elasticsearch và Kibana giữ nguyên volume.

\section{Kiểm thử tự động và bảo mật}

\begin{table}[H]
\centering
\caption{Kết quả kiểm tra release}
\label{tab:release-tests}
\begin{tabularx}{\textwidth}{Y C{4.0cm}}
\toprule
\textbf{Kiểm tra} & \textbf{Kết quả} \\
\midrule
Internal và public health & HTTP 200 \\
TLS public endpoint & Xác minh thành công \\
Login page & HTTP 200 \\
Members API khi chưa đăng nhập & HTTP 401 đúng auth boundary \\
Bộ test tự động & 96 passed \\
Elasticsearch integration CI & Passed \\
Trivy image scan & 0 High, 0 Critical \\
Kiến trúc image & linux/amd64 và linux/arm64 trên GHCR \\
Kích thước image & khoảng 146,2 MiB \\
\bottomrule
\end{tabularx}
\end{table}

\section{Pipeline thành công}

Run minh chứng: \url{https://github.com/AndyAnh174/dataops-demo-app/actions/runs/34995816200}.

\begin{table}[H]
\centering
\caption{Kết quả pipeline success}
\label{tab:success-run}
\begin{tabularx}{\textwidth}{Y C{5cm}}
\toprule
\textbf{Chỉ số} & \textbf{Kết quả} \\
\midrule
Tổng thời gian job & 9 phút 38 giây \\
DataOps Agent step & 6 phút 35 giây \\
Lifecycle event / Data Quality report & 2 / 1 \\
Log Elasticsearch & 813 dòng \\
Run status / incident mới & SUCCESS / 0 \\
Ứng dụng sau deploy & HTTP 200, khoảng 0,256 giây \\
\bottomrule
\end{tabularx}
\end{table}

Run này chứng minh luồng checkout, quality, build, scan, publish, deploy và telemetry hoạt động.
Happy path không gọi LLM. Tuy nhiên thời gian setup Node/cache và build/scan image còn lớn, cần
benchmark và tối ưu cache.

\section{Pipeline lỗi có kiểm soát}

Run minh chứng: \url{https://github.com/AndyAnh174/dataops-demo-app/actions/runs/34997172237}.
Kịch bản tiêm lỗi là \code{null\_rate} vượt ngưỡng ở stage Data Quality.

\begin{table}[H]
\centering
\caption{Kết quả pipeline failure}
\label{tab:failure-run}
\begin{tabularx}{\textwidth}{Y C{5cm}}
\toprule
\textbf{Chỉ số} & \textbf{Kết quả} \\
\midrule
Tổng thời gian job & 5 phút 8 giây \\
Thời gian Agent đến lúc phát hiện lỗi & khoảng 30 giây \\
Stage lỗi & data-quality \\
Event / report / log & 2 / 1 / 56 \\
Incident mới & 1 \\
Evidence & 4, không warning, không duplicate \\
Build/publish/deploy sau lỗi & Không chạy \\
Revision cũ & Vẫn healthy \\
\bottomrule
\end{tabularx}
\end{table}

Kết quả cho thấy quality gate đã ngăn release lỗi đi vào runtime. Evidence Collector thu đủ bốn
nhóm bằng chứng và BGE-M3 tạo knowledge document trong khoảng 60 giây.

\section{Đánh giá LLM và thử nghiệm function calling}

Hai model RCA được chạy trên cùng incident và cùng prompt/schema \code{rca-v1}:

\begin{table}[H]
\centering
\caption{So sánh model RCA trên incident thực}
\label{tab:models}
\begin{tabularx}{\textwidth}{L{3.2cm} C{2.8cm} C{1.8cm} Y}
\toprule
\textbf{Model} & \textbf{Thời gian} & \textbf{HTTP} & \textbf{Kết quả} \\
\midrule
\code{gemma4:e2b} & khoảng 176 giây & 422 & Output không vượt structured/schema validation \\
\code{qwen3.5:4b} & 29,23 giây & 422 & \code{LLM output failed the RCA schema validation} \\
\bottomrule
\end{tabularx}
\end{table}

Qwen nhanh hơn xấp xỉ sáu lần nhưng chưa chính xác hơn ở contract RCA. Runtime demo tạm giữ Qwen
để phục vụ vòng sửa prompt/schema; embedding vẫn dùng \code{bge-m3:567m}.

\code{functiongemma:latest} (khoảng 268 triệu tham số, Q8\_0) cũng được kiểm tra độc lập. Model
nhận biết được tool phù hợp trong một số prompt nhưng tạo duplicate tool call, có argument sai kiểu,
sinh thêm nội dung rác và dừng do giới hạn độ dài. Test single-tool vẫn lặp lời gọi. Vì vậy model
này chưa được tích hợp vào Control Plane. Kết luận kỹ thuật là capability ``tools'' trong metadata
không đồng nghĩa function calling đủ ổn định cho hệ thống có side effect.

% =============================================================================
\chapter{Đánh giá kết quả}

\section{Kết quả đạt được}

\begin{itemize}
  \item Kiến trúc tách Agent và Platform, core không khóa vào GitHub.
  \item Telemetry có correlation và idempotency; failure tạo incident đúng.
  \item Evidence có provenance, checksum, redact và giới hạn kích thước.
  \item Hybrid Retrieval kết hợp lexical/semantic thay vì đưa toàn bộ log vào context.
  \item Web UI có login, phân quyền, project/token, run/incident và recovery control.
  \item Luồng deploy dùng immutable image, backup, health gate và rollback artifact.
  \item Quality gate chặn release lỗi; revision đang chạy không bị ảnh hưởng.
  \item LLM bị cô lập khỏi credential và side effect; recovery có approval/audit/verification.
\end{itemize}

\section{Hạn chế và rủi ro}

\begin{enumerate}
  \item \textbf{RCA reliability:} schema validation thất bại với hai model; chưa có RCA report hợp
  lệ cho incident demo.
  \item \textbf{State khi phân tích lỗi:} incident còn ở \code{ANALYZING} sau HTTP 422, dễ gây
  hiểu nhầm tác vụ đang chạy.
  \item \textbf{Benchmark chưa đủ:} mới công bố đầy đủ một success và một \code{null\_rate}
  failure; chưa có bảng baseline trên F01--F06.
  \item \textbf{Portability chưa được chứng minh thực nghiệm:} provider-neutral ở mức thiết kế,
  nhưng chưa có GitLab/Jenkins adapter chạy thật.
  \item \textbf{Phân phối:} image đã có trên GHCR, job Docker Hub còn skip do thiếu biến/secret.
  \item \textbf{Hiệu năng CI:} restore cache trên self-hosted runner chiếm phần lớn thời gian ở ca
  lỗi nhanh; stage timestamp chưa đủ chi tiết.
  \item \textbf{Bảo mật vận hành:} credential từng được chia sẻ ngoài secret manager phải được
  thu hồi/rotate trước khi xem là môi trường production.
\end{enumerate}

\section{Đánh giá định lượng hiện tại}

\begin{table}[H]
\centering
\caption{Điểm đánh giá release demo}
\label{tab:scores}
\begin{tabularx}{\textwidth}{L{4.5cm} C{2.0cm} Y}
\toprule
\textbf{Hạng mục} & \textbf{Điểm} & \textbf{Nhận xét} \\
\midrule
Deployment safety & 9/10 & Backup, immutable tag, health gate, rollback rõ ràng \\
Telemetry/correlation & 9/10 & Event, report và log khớp đúng run \\
Failure containment & 9/10 & Pipeline lỗi không build/publish/deploy \\
Web/auth foundation & 8/10 & Nền tảng tốt; còn UAT đa tài khoản \\
AI RCA reliability & 4/10 & Evidence/RAG đạt, structured RCA thất bại \\
Performance & 6/10 & Setup/cache và image build/scan còn tốn thời gian \\
Operability & 7/10 & Dashboard, ELK, rollback có; thiếu trạng thái analysis failed \\
\midrule
\textbf{Tổng quan} & \textbf{7,4/10} & \textbf{Demo kỹ thuật tốt; AI RCA chưa production-ready} \\
\bottomrule
\end{tabularx}
\end{table}

% =============================================================================
\chapter{Kế hoạch công việc tiếp theo}

\section{Ưu tiên P0: hoàn thiện RCA an toàn}

\begin{enumerate}
  \item Ghi chi tiết lỗi Pydantic ở log nội bộ và trả failure category an toàn qua API.
  \item Thêm tối đa một lượt JSON repair với validation feedback; tuyệt đối không retry vô hạn.
  \item Thêm trạng thái \code{ANALYSIS\_FAILED}/\code{ACTION\_REQUIRED} khi cả lượt chính và repair
  đều thất bại.
  \item Bổ sung test regression bằng output thực của Gemma/Qwen: JSON sai, enum sai, citation hoặc
  knowledge ID không tồn tại.
  \item Chạy lại incident demo và chỉ đánh dấu M5 hoàn thành thực nghiệm khi RCA được lưu hợp lệ.
\end{enumerate}

\section{Ưu tiên P1: benchmark, UI và vận hành}

\begin{enumerate}
  \item Chạy F01--F06 lặp lại; đo root-cause accuracy, evidence recall@k, unsupported-claim rate,
  RCA latency, recovery success và unsafe action rate.
  \item So sánh rule-based, LLM raw log, LLM + Hybrid RAG và Agentic Hybrid RAG.
  \item UAT với OWNER/OPERATOR/VIEWER cho member, project deletion, token revoke và approval.
  \item Benchmark cache trên self-hosted runner; bổ sung \code{stage.started/completed} để đo timeline.
  \item Rotate credential, dùng quyền tối thiểu và chuyển runtime secret sang secret manager.
\end{enumerate}

\section{Ưu tiên P2: phân phối và đa nền tảng}

\begin{enumerate}
  \item Cấu hình Docker Hub access token và phát hành image SemVer đa kiến trúc.
  \item Hoàn thiện install/upgrade/backup/rollback guide và smoke test trên máy sạch.
  \item Triển khai GitLab CI hoặc Jenkins adapter, chạy lại cùng success/failure scenario.
  \item Đối chiếu mã nguồn core trước/sau provider thứ hai để chứng minh mức độ tái sử dụng.
\end{enumerate}

\section{Kết quả kỳ vọng cho lần báo cáo kế tiếp}

Lần báo cáo tiếp theo cần có ít nhất: một RCA hợp lệ có citation, một recovery được phê duyệt và
verification pass, bảng benchmark F01--F06, UAT ba role, image Docker Hub có tag bất biến và một
pipeline trên provider thứ hai. Khi đó đề tài mới đáp ứng đầy đủ tuyên bố ``đa nền tảng tích hợp
AI Agent'' ở mức thực nghiệm.

% =============================================================================
\chapter{Kết luận}

Dự án đã vượt qua giai đoạn thiết kế ý tưởng và có một hệ thống end-to-end chạy thật. Các khối
ingestion, incident, evidence, Hybrid RAG, policy/recovery, Web UI, authentication và vận hành
container đã được xây dựng và kiểm thử. Hai pipeline demo chứng minh DataOps Agent thu thập đúng
telemetry và quality gate bảo vệ revision đang chạy.

Kết quả đồng thời chỉ ra giới hạn nghiên cứu có giá trị: model nhỏ có thể nhanh hơn nhưng vẫn không
đảm bảo tuân thủ schema hoặc function calling. Vì vậy kiến trúc không đặt niềm tin vào LLM cho
side effect; mọi output phải được validate và mọi recovery phải qua deterministic policy,
approval và verification. Đây là đóng góp quan trọng về mặt kỹ thuật và an toàn của đề tài.

Ở thời điểm hiện tại, hệ thống phù hợp cho demo kỹ thuật và tiếp tục thực nghiệm. Công việc trọng
tâm không phải bổ sung thêm nhiều tính năng, mà là làm RCA đo lường được và fail-safe, hoàn thiện
benchmark, phát hành self-hosted ổn định và chứng minh provider thứ hai.

% =============================================================================
\appendix
\chapter{Minh chứng và đường dẫn}

\begin{longtable}{L{4.5cm} L{10.0cm}}
\toprule
\textbf{Nội dung} & \textbf{Đường dẫn/định danh} \\
\midrule
Repository chính & \url{https://github.com/AndyAnh174/dataops-control-plane} \\
DataOps Agent & \url{https://github.com/AndyAnh174/dataops-agent} \\
Web Console & \url{https://dataops-console.andyanh.id.vn} \\
Ứng dụng demo & \url{https://dataops.andyanh.id.vn} \\
Pipeline success & \url{https://github.com/AndyAnh174/dataops-demo-app/actions/runs/34995816200} \\
Pipeline failure & \url{https://github.com/AndyAnh174/dataops-demo-app/actions/runs/34997172237} \\
Release image & \code{ghcr.io/andyanh174/dataops-control-plane:sha-5575c6e31bdc62c65bbba3a07bd0bf62f9f24577} \\
License & MIT License, Copyright 2026 AndyAnh174 \\
\bottomrule
\end{longtable}

\chapter{API và cấu hình chính}

\begin{table}[H]
\centering
\caption{Nhóm endpoint quan trọng}
\begin{tabularx}{\textwidth}{L{5.8cm} Y}
\toprule
\textbf{Nhóm} & \textbf{Endpoint tiêu biểu} \\
\midrule
Event & \code{POST /api/v1/events/pipeline} \\
Log & \code{POST /api/v1/runs/\{run\_id\}/logs} \\
Data Quality & \code{POST /api/v1/runs/\{run\_id\}/reports/data-quality} \\
Incident/Evidence & \code{POST /api/v1/incidents/\{id\}/collect-evidence} \\
Retrieval & \code{POST /api/v1/retrieval/search} \\
RCA & \code{POST /api/v1/incidents/\{id\}/analyze} \\
Recovery & create plan, approve, execute, verification callback \\
Onboarding & \code{GET /api/v1/projects/\{id\}/onboarding/github} \\
\bottomrule
\end{tabularx}
\end{table}

Hai GitHub Secrets tối thiểu cho repository tích hợp là:

\begin{verbatim}
DATAOPS_URL=https://dataops-console.andyanh.id.vn
DATAOPS_TOKEN=<project-integration-token>
\end{verbatim}

Workflow gọi agent sau bước checkout:

\begin{verbatim}
- uses: AndyAnh174/dataops-agent@v0
  env:
    DATAOPS_URL: ${{ secrets.DATAOPS_URL }}
    DATAOPS_TOKEN: ${{ secrets.DATAOPS_TOKEN }}
\end{verbatim}

\chapter{Thuật ngữ viết tắt}

\begin{longtable}{L{3.0cm} L{11.0cm}}
\toprule
\textbf{Thuật ngữ} & \textbf{Ý nghĩa} \\
\midrule
CI/CD & Continuous Integration / Continuous Delivery \\
DataOps & Phương pháp tự động hóa và quản trị vòng đời pipeline dữ liệu \\
DQ & Data Quality \\
RCA & Root Cause Analysis -- phân tích nguyên nhân gốc \\
RAG & Retrieval-Augmented Generation \\
RRF & Reciprocal Rank Fusion \\
RBAC & Role-Based Access Control \\
MTTD & Mean Time To Detect \\
MTTR & Mean Time To Recover/Resolve \\
UAT & User Acceptance Testing \\
LLM & Large Language Model \\
\bottomrule
\end{longtable}

\begin{thebibliography}{9}
\bibitem{project-readme}
AndyAnh174, ``DataOps Control Plane -- README và tài liệu dự án,'' 2026.\\
\url{https://github.com/AndyAnh174/dataops-control-plane}

\bibitem{deployment-report}
AndyAnh174, ``Báo cáo triển khai và đánh giá DataOps end-to-end,'' tài liệu nội bộ dự án,
15--16/09/2026.

\bibitem{fastapi}
FastAPI, ``FastAPI Documentation.''\\
\url{https://fastapi.tiangolo.com/}

\bibitem{elastic}
Elastic, ``Elasticsearch Reference.''\\
\url{https://www.elastic.co/guide/en/elasticsearch/reference/current/index.html}

\bibitem{langgraph}
LangChain, ``LangGraph Documentation.''\\
\url{https://docs.langchain.com/oss/python/langgraph/}

\bibitem{ollama}
Ollama, ``Ollama Documentation.''\\
\url{https://docs.ollama.com/}
\end{thebibliography}

\end{document}
